\subsection{Aspectos Avanzados de Persistencia y Gestión de Archivos}
\label{subsec:aspectos_avanzados_persistencia}

La implementación completa del componente \textbf{FileManager} dentro del sistema de adquisición abarca una serie de mecanismos de ingeniería de software orientados a asegurar la robustez del almacenamiento frente a fallos y a optimizar el rendimiento de la base de datos de ensayos. A continuación, se desarrollan detalladamente tres tópicos críticos del módulo que no fueron abordados en la descripción preliminar: la estructura jerárquica del sistema de ficheros en disco, la robustez del flujo de guardado de datos bajo condiciones adversas de interferencia electromagnética, y los métodos de recuperación y exportación de trazas para análisis externos.

\subsubsection{Estructura Jerárquica del Sistema de Ficheros en Disco}
\label{subsubsec:estructura_ficheros_disco}

El módulo \textbf{FileManager} automatiza la creación física de una estructura de directorios estandarizada a partir de un directorio raíz seleccionado por el operador para el proyecto de ensayo. Esta topología garantiza una segregación física limpia de los archivos generados y previene la dispersión involuntaria de la información. El árbol de directorios creado de manera determinista en el disco rígido posee el siguiente diseño:

\begin{itemize}
    \item \texttt{01\_Base\_Datos/}: Contiene el archivo principal con formato estructurado \textbf{HDF5} del proyecto (con extensión \texttt{.h5}), el cual unifica la metadata global de la sesión, los atributos de cada descarga y las grandes matrices de datos discretos.
    \item \texttt{02\_Respaldo/}: Aloja los archivos de respaldo temporal con codificación de texto plano (\textbf{CSV}). Por cada disparo registrado, se genera un archivo independiente para el canal de tensión (CH1) y corriente (CH2) conteniendo la traza digitalizada cruda obtenida directamente de la memoria del osciloscopio digital \cite{gds1000au_ProgrammingManual}.
    \item \texttt{03\_Resultados/}: Almacena los cuadros de datos consolidados exportados de manera manual por el operador en formato de hojas de cálculo tipo \textbf{XLSX} o archivos tabulares compatibles de texto plano \textbf{CSV} para su procesamiento complementario en otras herramientas analíticas.
    \item \texttt{04\_Reportes/}: Carpeta destinada al almacenamiento de los informes técnicos de validación y calibración de software metrológico automatizados en formato portable de documento de lectura (\textbf{PDF}), generados a partir de los resultados calculados por el sistema \cite{iec61083_2}.
\end{itemize}

El método encargado de esta tarea es \texttt{create\_project\_directories(base\_path)}. Este método encapsula llamadas de bajo nivel al sistema operativo mediante la clase \textbf{Path} del paquete nativo \textbf{pathlib}, garantizando que, si las carpetas ya existen, no se interrumpa el flujo del programa principal ni se altere el contenido preexistente.

\subsubsection{Robustez del Almacenamiento e Integridad frente a Fallos de Redundancia}
\label{subsubsec:robustez_integridad_fallos}

En el entorno físico de un laboratorio de alta tensión, la ocurrencia de descargas disruptivas tipo rayo induce transitorios electromagnéticos de gran energía (interferencia electromagnética o \textbf{EMI}). Estos fenómenos pueden comprometer la estabilidad del sistema informático, provocando desde caídas de tensión hasta el congelamiento del software de host. 

Para proteger la integridad física y lógica de los registros, el módulo \textbf{FileManager} incorpora un flujo atómico de guardado estructurado mediante el uso de manejadores de contexto seguros (\texttt{with}) provistos por la biblioteca \textbf{h5py}. La secuencia operativa para persistir cada onda mitigando riesgos de corrupción se define formalmente en los siguientes pasos lógicos:

\begin{enumerate}
    \item \textbf{Apertura Segura}: El archivo HDF5 se abre dinámicamente utilizando el modo de append (\texttt{'a'}). El manejador de archivos de Python asegura la liberación del descriptor de archivos en el sistema operativo incluso si se produce un fallo crítico o una excepción durante la escritura, mitigando el bloqueo definitivo de la base de datos.
    \item \textbf{Verificación y Sanitización}: La metadata tipeada por el operador se valida sintácticamente y se filtra de forma interactiva empleando expresiones regulares \textbf{re} antes de su inserción binaria.
    \item \textbf{Inyección de Atributos}: Se crea un nuevo grupo interno para el disparo (p. ej. \texttt{/Onda\_001} o \texttt{/Referencia\_001}). Los parámetros ambientales, condiciones eléctricas de hardware y la marca de tiempo se inyectan como atributos de este grupo específico, garantizando que cada traza digitalizada esté indisolublemente asociada a sus condiciones físicas de medición reales \cite{iec60060_1}.
    \item \textbf{Optimización de Memoria (Chunking)}: Al crear los conjuntos de datos lógicos (\textbf{datasets}) para almacenar los vectores densos NumPy de tensión y corriente, el módulo configura dinámicamente un esquema de segmentación (\textbf{chunking}) y una compresión sin pérdida tipo \textbf{GZIP} \cite{file_manager_txt}:
    \begin{equation}
    \label{eq:chunk_size}
    \text{Chunk\_Size} = N_{record} \cdot \qty{2}{\text{bytes}}
    \end{equation}
    Donde $N_{record}$ representa la longitud de registro configurada para el hardware. Este esquema optimiza drásticamente la velocidad de almacenamiento en disco para matrices densas de hasta \num{2000000} de muestras de canal único \cite{gds1000au_Spec}.
    \item \textbf{Escritura Redundante y Sincronización}: Al finalizar la inyección de cada canal analógico activo, el módulo realiza un volcado físico inmediato (\texttt{flush()}) de los buffers de memoria del disco físico, de manera que la información se guarde de forma permanente, seguido por la escritura simultánea del archivo CSV de respaldo en el directorio \texttt{02\_Respaldo/}.
\end{enumerate}

Cualquier excepción de entrada y salida, como la falta de espacio en disco o el bloqueo por escritura redundante de HDF5, se captura mediante un bloque estructurado \texttt{try-except OSError}, redirigiendo el suceso a la interfaz de usuario en forma de un cuadro de advertencia visual sin interrumpir la ejecución del programa principal ni del sub-hilo secundario de polling del osciloscopio.

\subsubsection{Métodos de Recuperación y Exportación Tabular de Trazas}
\label{subsubsec:recuperacion_exportacion_trazas}

El análisis metrológico histórico y la validación cruzada algorítmica exigen subrutinas de alta eficiencia para recuperar las trazas almacenadas en formato binario. Para resolver este requerimiento metrológico, el módulo \textbf{FileManager} provee métodos de lectura indexada que actúan de manera inversa al proceso de guardado.

El método \texttt{get\_waveform\_data(project\_path, group\_name)} accede directamente al archivo HDF5 en modo lectura exclusiva (\texttt{'r'}), el cual previene conflictos de concurrencia de lectura con otros hilos de ejecución. El método navega jerárquicamente por la estructura interna del archivo conforme al mapa lógico estático \texttt{HDF5\_DATASET\_MAP} y devuelve los vectores numéricos crudos y procesados en forma de matrices estructuradas NumPy.

Para la exportación de resultados a hojas de cálculo comerciales tipo Excel (\texttt{.xlsx}) o de texto plano (\texttt{.csv}), el sistema recupera síncronamente todas las descargas del proyecto aplicando factores de atenuación eléctrica de los divisores resistivos e instrumental coaxial físico externos. El escalamiento final de magnitudes temporales y eléctricas se expresa formalmente como:

\begin{equation}
\label{eq:escalado_final_export}
u_{export}(t) = u_{raw}(t) \cdot F_{div} \cdot F_{att} \cdot \qty{10^{-3}}{\text{kV/V}}
\end{equation}

Donde las variables implicadas representan:
\begin{itemize}
    \item $u_{export}(t)$: Vector final de tensión escalada expresada en kilovoltios (\unit{\kilo\volt}).
    \item $u_{raw}(t)$: Tensión instantánea digitalizada por el osciloscopio en voltios (\unit{\volt}).
    \item $F_{div}$: Factor de escala de atenuación física del divisor de tensión resistivo, fijado en \num{614}.
    \item $F_{att}$: Factor de atenuación analógica de hardware, fijado nominalmente en \num{50}.
\end{itemize}

Durante esta fase, la biblioteca \textbf{pandas} encapsula el ensamblado de las tablas dinámicas, indexando de manera cronológica cada disparo y guardando los resultados consolidados con soporte de la librería de bajo nivel \textbf{openpyxl}. Este flujo desacoplado asegura que los datos crudos históricos inmutables almacenados en HDF5 permanezcan completamente inalterados frente a cualquier manipulación de escalado para visualización o exportación.

\begin{thebibliography}{99}
    \bibitem{iec60060_1} International Electrotechnical Commission, ``High-voltage test techniques - Part 1: General definitions and test requirements'', IEC 60060-1, Ed. 3.0, 2010.
    \bibitem{iec61083_2} International Electrotechnical Commission, ``Instruments and software used for measurement in high-voltage impulse tests - Part 2: Evaluation of software used for the determination of the parameters of impulse waveforms'', IEC 61083-2, Ed. 2.0, 2013.
    \bibitem{gds1000au_ProgrammingManual} Good Will Instrument Co., Ltd., ``GDS-1000A-U Series Digital Storage Oscilloscope Programming Manual'', Part No. 82DS-112AUI01, Taipei, Taiwan, 2011.
    \bibitem{gds1000au_Spec} Good Will Instrument Co., Ltd., ``GDS-1000A-U Specifications Sheet'', Taipei, Taiwan, 2011.
    \bibitem{file_manager_txt} Excerpts from ``file_manager.txt'': Módulo de almacenamiento y administración persistente del sistema de ficheros del ensayo, 2026.
\end{thebibliography}